\documentclass[11pt,letterpaper]{article}
\usepackage[T1]{fontenc}
\usepackage{lmodern,amsmath,amssymb,amsthm,microtype}
\usepackage[margin=1in]{geometry}
\usepackage[hidelinks]{hyperref}
\newtheorem{theorem}{Theorem}
\newtheorem{lemma}{Lemma}
\title{A fractional improvement for nine separated subset sums}
\author{Aafi Khan Malek\\ \small Independent researcher\\
\small\href{mailto:aafimalek2032@gmail.com}{\texttt{aafimalek2032@gmail.com}}}
\date{27 November 2026}
\begin{document}
\maketitle
\begin{abstract}
Let $f_{\mathbb R}(9)$ be the smallest possible maximum of nine
positive real weights whose subset sums are pairwise at least one apart.
We show that $f_{\mathbb R}(9)\ge148+1/6$, improving the bandwidth
bound of $148$. A block of consecutive subset sums forced by hypercube
bandwidth constrains all nine fractional parts. A second-moment argument
then bounds the number of terms in a collision after rounding down.
The proof uses the published integer optimum $161$ and does not settle
the proposed bound $149$.
\end{abstract}

\begin{theorem}
If $0<x_1<\cdots<x_9$ and $|s(A)-s(B)|\ge1$ for distinct
$A,B\subseteq\{1,\ldots,9\}$, where $s(A)=\sum_{i\in A}x_i$,
then $x_9\ge148+1/6$.
\end{theorem}

\begin{proof}
The bandwidth of the nine-dimensional hypercube is
$\sum_{j=0}^8\binom{j}{\lfloor j/2\rfloor}=148$
\cite{harper,costa}. Label its vertices by the ranks (in increasing
order) of their subset sums. If an edge changes coordinate $i$, its
rank difference is at most $x_i$, since consecutive sums are at least
one apart. In particular $M:=x_9\ge148$. We may assume $M=148+\delta$
with $0\le\delta<1/6$.

\paragraph{A critical block.}
Separation of singleton sums gives $x_i\le M-1<148$ for $i\le8$.
Thus an edge with rank difference at least $148$ must change coordinate
$9$, and the rank difference is exactly $148$. The closed interval
between its two sums contains exactly $149$ consecutive sums.
For two sums in this interval with rank difference $r$, their
difference lies in $[r,r+\delta]$: the total excess of the $148$
consecutive gaps over one is exactly $\delta$.

Delete coordinate $9$ from each of the $149$ vertices in this block.
Only the endpoints can have the same projection, because two vertices
with equal projections differ in sum by exactly $M$, the entire
length of the interval. The resulting $148$ distinct vertices occupy
more than half of the eight-dimensional cube. In each of the eight
coordinate directions, the matching of the cube along that coordinate
therefore contains an edge whose two vertices are in the projection.
Choose representatives for its endpoints in the block. Their positive
difference is $x_i$ or $M-x_i$, since $M+x_i$ cannot fit inside the
interval. Every difference in the block has fractional part at most
$\delta$. If the difference is $M-x_i$ and the fractional part
$\theta_i$ of $x_i$ were greater than $\delta$, the fractional part
of $M-x_i$ would be $1+\delta-\theta_i>\delta$. Hence
$0\le\theta_i\le\delta$ for each $i$; this also holds for $i=9$.

\paragraph{A second-moment restriction.}
For the $N=512$ sorted sums $t_1,\ldots,t_N$, spacing gives
$t_j-t_i\ge j-i$ for $j>i$. By the pairwise formula for variance,
\[
 \frac{1}{N}\sum_{j=1}^N (t_j-\bar t)^2
 =\frac{1}{N^2}\sum_{i<j}(t_j-t_i)^2
 \ge\frac{N^2-1}{12}.
\]
A uniform subset includes each coordinate independently with
probability $1/2$, so the same variance is $\frac14\sum_{i=1}^9x_i^2$.
It follows that $\sum_i x_i^2\ge87381$.

Put $a_i=\lfloor x_i\rfloor$. These are positive integers since even
each singleton sum is at least one from the empty sum. Suppose two
subsets have the same sum of $a_i$. Cancel their intersection. The two
remaining subsets are disjoint and nonempty. A side of size eight
would have real sum at least $255$, since its $256$ subset sums run
from $0$ to its total and are separated by at least one. But equality
of the rounded sums with a singleton bounds that total by
$148+8\delta<156$, a contradiction.

For a side of size seven the other side has at most two entries.
Hence the seven real weights have total less than
$2\cdot148+7=303$. Write these weights in decreasing order as
$z_1,\ldots,z_7$. The last $7-k$ of them have sum at least
$2^{7-k}-1$ by separation of their own subset sums. Since
$z_j<149$, the first $k$ have sum bounded above by
\[
 \min\{149k,303-(2^{7-k}-1)\}.
\]
For $k=1,\ldots,7$ these are the prefix sums of
$(149,123,16,8,4,2,1)$. If nonnegative decreasing vectors $z,w$
satisfy these prefix inequalities, then for all $t\ge0$,
\[
 \sum_j(z_j-t)_+=\max_{0\le k\le7}\Big(\sum_{j\le k}z_j-kt\Big)
 \le\sum_j(w_j-t)_+.
\]
Integration over $t\ge0$ gives $\sum_j z_j^2\le\sum_j w_j^2=37671$.
Even allowing $149^2$ for each remaining squared weight gives
$\sum_i x_i^2\le82073$, contradicting the second-moment restriction.

\paragraph{Rounding.}
Each side of a rounded collision must therefore have size at most
six. Its real sum differs from its rounded sum by a number in
$[0,6\delta]$. The two real sums would differ by at most
$6\delta<1$, contradicting separation. Consequently the $a_i$ form
a nine-element integer set with distinct subset sums and maximum
$148$. This contradicts the exact integer optimum $161$
\cite[Table 5]{bm}.
\end{proof}

\paragraph{Scope and reproducibility.}
The critical-block argument also gives a general observation. Let $H_n$
be the bandwidth of the $n$-cube. If $H_n>2^{n-2}$ and the largest
weight is $H_n+\delta$ with $0\le\delta<1$, projection of its critical
block again covers more than half the $(n-1)$-cube. Thus every
fractional part is at most $\delta$. If the integer optimum exceeds
$H_n$, rounding down gives a contradiction for $\delta<1/(n-1)$:
after cancellation either side of a putative collision has at most
$n-1$ terms. The second-moment step in the theorem sharpens this
general estimate at $n=9$.

The companion script verifies the numerical constants and enumerates
the $512$ sums of the classical upper-bound example. It does not
independently reprove the hypercube bandwidth formula or the published
integer optimum. The proposed bound $149$ is discussed in
\cite[Section 7]{notik}. This note is not peer reviewed or formalized
in Lean.

\paragraph{Contributions and AI disclosure.}
Aafi Khan Malek selected and directed the research project and requested
the repository and verification workflow. ChatGPT assisted with the
literature search, argument, write-up, and verification script. Codex
was not used.

\begin{thebibliography}{9}
\bibitem{harper} L. H. Harper, Optimal numberings and isoperimetric
problems on graphs, J. Combinatorial Theory \textbf{1} (1966), 385--393.
\bibitem{costa} S. Costa and S. Della Fiore, A bandwidth refinement of
the Erd\H{o}s distinct subset sums bound,
\href{https://arxiv.org/abs/2609.27941}{arXiv:2609.27941} (2026).
\bibitem{bm} P. Borwein and M. J. Mossinghoff, Newman polynomials with
prescribed vanishing and integer sets with distinct subset sums,
\href{https://www.cecm.sfu.ca/~pborwein/PAPERS/P170.pdf}{Math. Comp.
\textbf{72} (2003), 787--800}.
\bibitem{notik} E. Notik, Positive dual certificates for real separated
subset sums: chamber optimality of the Conway--Guy and Bohman sets,
and exact values up to eight,
\href{https://arxiv.org/abs/2611.17360}{arXiv:2611.17360} (2026).
\end{thebibliography}
\end{document}
